\label{chap:planck-autodual}

\section{Réalisation de Compton associée à la périodicité temporelle de \(108\) phases}

Sur le cycle de \cref{eq:realizacion-circular-ct108}, nous réutilisons la section
positive d’action \(\hbar\) déjà engendrée par la branche déterminantale d’action.
Le quantum angulaire associé est
\begin{equation}
E_{108}=\hbar\omega_{108},
\qquad
m_{108}=\frac{E_{108}}{c^2}
=\frac{\hbar\omega_{108}}{c^2}.
\label{eq:cuanto-ct108}
\end{equation}
\(G\) n’a pas encore été utilisé. La géométrie du cycle et la quantification
angulaire suffisent à fixer ses deux lecteurs de Compton.

\begin{theorem}[Identité entre la périodicité temporelle de \(108\) phases et la longueur de Compton]
\label{thm:ct108-compton}
Pour le quantum \eqref{eq:cuanto-ct108},
\begin{equation}
\lambdabarC(m_{108})=\frac{\hbar}{m_{108}c}=R_{108},
\qquad
\lambdaC(m_{108})=\frac{h}{m_{108}c}=C_{108}.
\label{eq:identidad-ct108-compton}
\end{equation}
Par conséquent, le rayon du cycle est sa longueur de Compton réduite et la
longueur d’un tour est sa longueur de Compton complète.
\end{theorem}

\begin{proof}
En substituant \(m_{108}=\hbar\omega_{108}/c^2\),
\[
\frac{\hbar}{m_{108}c}
=\frac{\hbar c^2}{\hbar\omega_{108}c}
=\frac{c}{\omega_{108}}=R_{108}.
\]
Comme \(h=2\pi\hbar\), la deuxième égalité est \(2\pi R_{108}=C_{108}\).
\end{proof}

La coïncidence n’est pas un ajustement numérique. C’est une identité homogène : elle vaut
pour tout \(\hbar,c,t_0>0\) et dépend seulement du fait que la même période détermine
la fréquence, la longueur du cycle et le quantum d’action.

\section{Transducteur gravitationnel et équation de sélection autoduale}

La droite gravitationnelle HMT--MD contient le transducteur
\begin{equation}
G(\theta)=\theta\frac{c^5t_0^2}{\hbar},
\qquad \theta>0.
\label{eq:transductor-g-planck-anexo}
\end{equation}
L’homogénéité est exacte, mais ne sélectionne pas à elle seule \(\theta\). Le
cycle CT108 apporte une condition interne supplémentaire : le lecteur gravitationnel
doit renvoyer la même échelle radiale que le lecteur géométrique et le lecteur de
Compton.

Nous définissons le rayon gravitationnel réduit
\begin{equation}
\rg(\theta):=\frac{G(\theta)m_{108}}{c^2}.
\label{eq:radio-gravitatorio-reducido}
\end{equation}
En substituant \eqref{eq:cuanto-ct108} et
\eqref{eq:transductor-g-planck-anexo},
\begin{equation}
\rg(\theta)
=\theta\frac{\pi}{54}\ell_0.
\label{eq:rg-theta-lineal}
\end{equation}

\begin{theorem}[Sélecteur autodual de la périodicité temporelle de \(108\) phases]
\label{thm:selector-autodual-ct108}
L’équation
\[
\rg(\theta)=R_{108}
\]
possède une unique solution positive :
\begin{equation}
\boxed{\thetaStar=\left(\frac{54}{\pi}\right)^2.}
\label{eq:theta-star}
\end{equation}
\end{theorem}

\begin{proof}
Par \eqref{eq:realizacion-circular-ct108} et
\eqref{eq:rg-theta-lineal}, l’équation est
\[
\theta\frac{\pi}{54}\ell_0
=\frac{54}{\pi}\ell_0.
\]
Comme \(\ell_0>0\), on simplifie et l’on obtient
\(\theta=(54/\pi)^2\). La fonction
\(\theta\mapsto\rg(\theta)\) est strictement croissante sur
\(\RR_{>0}\), de sorte qu’il n’existe aucune autre solution positive.
\end{proof}

\begin{corollary}[Unité de Planck interne]
\label{cor:unidad-planck-interna}
Avec \(G_*=G(\thetaStar)\), les sections de Planck satisfont
\begin{equation}
\ellP:=\sqrt{\frac{\hbar G_*}{c^3}}=R_{108},
\qquad
\tP:=\sqrt{\frac{\hbar G_*}{c^5}}=\frac1{\omega_{108}},
\label{eq:planck-long-time-ct108}
\end{equation}
\begin{equation}
\mP:=\sqrt{\frac{\hbar c}{G_*}}=m_{108},
\qquad
\EP:=\mP c^2=E_{108}.
\label{eq:planck-mass-energy-ct108}
\end{equation}
En particulier,
\begin{equation}
\boxed{
R_{108}=\lambdabarC=\rg=\ellP,
\qquad
C_{108}=\lambdaC=2\pi\rg=2\pi\ellP.}
\label{eq:tres-lectores-planck}
\end{equation}
\end{corollary}

\begin{proof}
On substitue \(G_*=(54/\pi)^2c^5t_0^2/\hbar\). Alors
\[
\ellP=\frac{54}{\pi}ct_0=R_{108},
\qquad
\tP=\frac{54}{\pi}t_0=\frac1{\omega_{108}},
\]
et
\[
\mP=\frac{\pi}{54}\frac{\hbar}{c^2t_0}
=\frac{\hbar\omega_{108}}{c^2}=m_{108}.
\]
Le reste découle du \cref{thm:ct108-compton}.
\end{proof}

\begin{remark}[Trois lecteurs, une section]
La géométrie circulaire publie \(R_{108}\) et \(C_{108}\) ; la quantification de
Compton publie \(\lambdabarC\) et \(\lambdaC\) ; la réalisation gravitationnelle
publie \(\rg\) et \(2\pi\rg\). En \(\thetaStar\), ces trois paires sont les
coordonnées d’une unique section. C’est la raison précise pour laquelle
l’unité de Planck cesse d’apparaître comme une combinaison dimensionnelle accidentelle
dans cette formalisation.
\end{remark}

\section{Équivalence des normalisations sur l’identité autoduale}

Le transducteur \eqref{eq:transductor-g-planck-anexo} admet deux manières de
fixer l’échelle élémentaire, qui ne doivent pas être confondues.

\begin{enumerate}[label=\textup{(\roman*)}]
\item Si l’on fixe \(\theta=1\), alors
\(\sqrt{\hbar G/c^3}=\ell_0\). La transition élémentaire \(ct_0\) est identifiée
à la longueur de Planck de cette carte.
\item Si l’on exige que l’\emph{orbite complète} CT108 ait le rayon de Planck,
la condition sélectionne \(\thetaStar=(54/\pi)^2\) et
\(R_{108}=\ellP\).
\end{enumerate}

Ce ne sont pas des valeurs rivales de \(G\). Ce sont deux jauges de la même équation
homogène : l’une attribue l’unité de Planck au pas élémentaire ; l’autre, au rayon
du cycle complet. L’annexe utilise la seconde parce que sa thèse concerne
l’identité de la clôture CT108.

\section{Conventions gravitationnelles : rayon réduit et rayon de Schwarzschild}

La notation gravitationnelle est fixée sans équivoque :
\[
\rg=\frac{Gm}{c^2},
\qquad
\rs=\frac{2Gm}{c^2}=2\rg.
\]
Le théorème précédent emploie \(\rg\), non \(\rs\). Dans la convention habituelle de
la masse de Planck, l’égalité \(\hbar/(mc)=Gm/c^2\) se produit en \(m=\mP\),
tandis que
\begin{equation}
\frac{\hbar}{mc}=\frac{2Gm}{c^2}
\quad\Longleftrightarrow\quad
m=\frac{\mP}{\sqrt2}.
\label{eq:cruce-schwarzschild-factor-dos}
\end{equation}
C’est pourquoi l’expression « échelle gravitationnelle » désigne dans la thèse directrice le rayon
réduit. Le rayon d’horizon conventionnel demeure un lecteur distinct.

\section{Dédoublement bidirectionnel de la droite d’action}

La droite interne d’action possède deux sections orientées :
\begin{equation}
\hbarpre=H_5\,10^{-34}\mathcal U_S,
\qquad
\hbarret=\eta_{\mathrm{ret}}\,10^{-34}\mathcal U_S,
\qquad
\Ract=\frac{H_5}{\eta_{\mathrm{ret}}}.
\label{eq:hbar-pre-ret-anexo}
\end{equation}
La décade \(10^{-34}\) fixe l’échelle absolue ; le quotient \(\Ract\)
transporte la différence orientée. Cette séparation permet de demander si
l’autodualité de Planck subsiste dans les deux coordonnées.

\begin{theorem}[Covariance réciproque de l’action et de la gravité]
\label{thm:accion-gravedad-reciproca}
Pour \(a\in\{\mathrm{pre},\mathrm{ret}\}\), soit
\[
G_a(\theta)=\theta\frac{c^5t_0^2}{\hbar_a},
\qquad
m_{108,a}=\frac{\hbar_a\omega_{108}}{c^2}.
\]
Alors
\begin{equation}
\frac{G_a(\theta)m_{108,a}}{c^2}
=\theta\frac{\pi}{54}\ell_0
\label{eq:invariante-pre-ret-planck}
\end{equation}
est indépendant de \(a\). De plus,
\begin{equation}
\frac{m_{108,\mathrm{pre}}}{m_{108,\mathrm{ret}}}=\Ract,
\qquad
\frac{G_{\mathrm{pre}}}{G_{\mathrm{ret}}}=\Ract^{-1},
\label{eq:reciprocidad-pre-ret}
\end{equation}
et la longueur de Planck
\(\sqrt{\hbar_aG_a/c^3}=\sqrt\theta\,\ell_0\) est la même dans les deux
orientations.
\end{theorem}

\begin{proof}
Les facteurs \(\hbar_a\) s’annulent dans le produit \(G_am_{108,a}\).
Les deux rapports de \eqref{eq:reciprocidad-pre-ret} sont obtenus directement
à partir des définitions. L’invariance de la longueur de Planck découle de
\(\hbar_aG_a=\theta c^5t_0^2\).
\end{proof}

Une paramétrisation symétrique rend la géométrie visible :
\begin{equation}
\hbar_{\mathrm{geom}}:=\sqrt{\hbarpre\hbarret},
\qquad
\xi:=\frac12\log\Ract,
\qquad
\hbarpre=\hbar_{\mathrm{geom}}e^\xi,
\quad
\hbarret=\hbar_{\mathrm{geom}}e^{-\xi}.
\label{eq:rapidez-accion-pre-ret}
\end{equation}
La mémoire bidirectionnelle réside dans \(\xi\) ; la géométrie de Planck
réside dans le produit invariant. Cela permet une différence réelle entre les
deux coordonnées d’action sans exiger une différence de vitesse de
propagation.

\begin{remark}[Portée physique du théorème]
La réciprocité de \cref{thm:accion-gravedad-reciproca} est exacte au sein du
transducteur. Pour identifier \(G_{\mathrm{pre}}\) et \(G_{\mathrm{ret}}\) à deux
mesures gravitationnelles orientées, il faut spécifier le protocole,
l’observable et la symétrie qui les relie. Le théorème prouve la
compatibilité structurelle ; il n’anticipe pas à lui seul une violation de Lorentz
ni une variation observable de \(G\).
\end{remark}
